\documentclass[11pt,a4paper]{article}
\usepackage[utf8]{inputenc}
\usepackage[T1]{fontenc}
\usepackage[italian]{babel}
\usepackage{amsmath,amssymb,amsthm,geometry,graphicx,hyperref}
\geometry{margin=2.5cm}
\newtheorem{proposizione}{Proposizione}
\title{MGD Neuro 0.33.0: memoria esperienziale multimodale\\Garanzie condizionate e limiti verificati}
\author{Nota tecnica sui sorgenti e sugli esperimenti}
\date{26 settembre 2026}
\begin{document}
\maketitle
\noindent Questa nota descrive il nuovo classificatore episodico integrato nell'app.
Non costituisce una dimostrazione dei paper originali MGD, di superiorità generale
su Transformer, di intelligenza generale o di efficienza energetica biologica.

\section{Esperienze, contesti e geometria}
Si conserva una successione finita $E=(e_1,\ldots,e_n)$ di esperienze confermate,
con $e_i=(c_i,y_i,\{z_i^{(m)}\}_{m\in M_i},p_i)$: contesto, etichetta,
descrittori normalizzati delle modalità presenti e provenienza.
Le etichette non vengono inserite automaticamente nei descrittori.
Con una metrica diagonale positiva $w_{cmk}\in[0.05,8]$, si usa la dissimilarità
\[
 D_{cm}(z,z_i)=\sum_k w_{cmk}(z_k-z_{ik})^2.
\]
La radice di questa espressione è una distanza sullo spazio dei descrittori;
la sua versione quadratica è quella usata nei kernel. La rappresentazione di
immagini o suoni può essere non iniettiva: distanza nulla non dimostra identità
degli oggetti nel mondo.

Per una query con modalità $M$, i confronti congiunti ammettono soltanto episodi
che contengono tutte le modalità richieste nello stesso contesto. Posto
$E_{cyM}=\{i:c_i=c,y_i=y,M\subseteq M_i\}$, si calcola
\[
 s_y(z)=\frac1{|E_{cyM}|}\sum_{i\in E_{cyM}}
 \exp\!\left[-\frac{\beta_c}{|M|}\sum_{m\in M}D_{cm}(z^{(m)},z_i^{(m)})\right],
 \qquad q_y(z)=\frac{s_y(z)}{\sum_{y'}s_{y'}(z)}.
\]
Si tratta di punteggi normalizzati su categorie osservate, non di probabilità
calibrate della verità nel mondo aperto. In assenza di esempi confrontabili il
modello si astiene. Si calcolano anche punteggi dei singoli canali per rendere
visibile il loro contributo. Mediare prodotti all'interno degli episodi conserva
associazioni che il prodotto di medie marginali perderebbe, come mostra il test XOR.

\section{Aggiornamento e verifica sui ricordi}
Per ciascun canale, l'aggiornamento propone un passo sulla perdita logaritmica
del kernel bilanciato per classe, su un insieme di al massimo 48 riferimenti.
Con $\delta_{ik}=(z_k-z_{ik})^2$, il gradiente rispetto a $w_k$ è
\[
 \frac{\partial(-\log q_y)}{\partial w_k}
 =\beta_c\bigl(\mathbb E[\delta_{ik}\mid y,z]
                  -\mathbb E[\delta_{ik}\mid z]\bigr).
\]
Si proietta poi il passo nell'intervallo consentito. Il candidato viene accettato
solo se non peggiora la perdita complessiva e non cambia in errata una decisione
prima corretta su al massimo 24 riferimenti, escludendo ogni riferimento dal suo
insieme di confronto. Questo controllo campionato non è una garanzia sulle
esperienze non controllate né sui dati futuri. Il passo ottimizza un obiettivo
per canale; la verifica usa anche la previsione congiunta.

Ogni otto esempi di un contesto, quando ogni categoria osservata dispone di
almeno due esempi, $\beta_c$ viene selezionato tra $6,12,24,48$ mediante perdita
leave-one-out su al massimo 16 riferimenti. Non si accede ai dati di verifica.
La selezione impedisce di trasformare in errata una previsione corretta dei
riferimenti controllati e penalizza debolmente valori elevati di $\beta_c$.

\section{Che cosa si può dimostrare}
\begin{proposizione}[Richiamo condizionato]
Se una query ha descrittori identici a quelli di almeno un episodio conservato,
nel contesto richiesto e su tutte le modalità della query, e tutte le
corrispondenti etichette conservate sono uguali a $y$, la regola di richiamo
esatto restituisce $y$, indipendentemente dagli aggiornamenti della metrica.
\end{proposizione}
\begin{proof}
Il codice verifica i descrittori prima della decisione per kernel. L'insieme
delle etichette esattamente corrispondenti è il singoletto $\{y\}$ e il ramo
di richiamo restituisce tale elemento, senza consultare i pesi. La tolleranza
numerica implementata è $10^{-9}$ per coordinata. Episodi con etichette
incompatibili producono invece un'astensione. La conclusione richiede che
l'episodio resti conservato e che non venga aggiunta una conferma incompatibile.
\end{proof}
La proposizione non garantisce riconoscimento di nuove viste di uno stesso
oggetto, assenza di collisioni del descrittore o assenza di dimenticanza
nella generalizzazione. Il limite di 2.048 episodi rifiuta ulteriori inserimenti;
non promette memoria illimitata.

\begin{proposizione}[Ricostruzione dopo eliminazione]
Fissati descrittori, ordine e algoritmo deterministico di aggiornamento,
azzerare lo stato adattivo e riprodurre gli episodi superstiti produce lo stesso
predittore che si sarebbe ottenuto addestrando da zero sui soli superstiti.
\end{proposizione}
\begin{proof}
Per induzione sulla lunghezza della successione superstite: gli stati iniziali
sono uguali. Al passo successivo sono uguali riferimenti, contesto, parametri,
proposta di aggiornamento e verifica; quindi anche lo stato successivo è uguale.
Gli identificatori assoluti e le date non partecipano alle distanze, al
campionamento ordinato o al criterio di accettazione. La selezione della
temperatura viene ripetuta agli stessi conteggi per contesto. I contatori di
valutazione ripartono da zero dopo l'eliminazione, evitando di presentarli come
misure originali di una storia che è stata modificata.
\end{proof}
La conclusione riguarda questo componente; non prova unlearning forense degli
altri moduli dell'app, di copie precedenti o del sistema operativo.

\begin{proposizione}[Limite della conservazione universale]
Un sistema con al più $b$ bit di stato persistente non può conservare esattamente
tutte le successioni binarie di lunghezza $n>b$ e rispondere senza errori a
qualunque futura domanda sul valore di un loro bit.
\end{proposizione}
\begin{proof}
Ci sono $2^n$ successioni e al più $2^b$ stati. Due successioni differenti hanno
quindi lo stesso stato. La domanda relativa a una posizione in cui differiscono
richiede due risposte diverse dallo stesso stato: almeno una sarà errata.
\end{proof}
Una compressione utile deve dunque specificare quali predizioni conservare,
con quale errore tollerato e rispetto a quale distribuzione. Non è stato
dimostrato né implementato un criterio universale di compressione in questa
versione. I vincoli di apprendimento senza fine, memoria finita e richiamo
perfetto di ogni dettaglio non possono essere promessi simultaneamente.

\section{Esperimento e risultato negativo sul primato}
Su sei seed finali, con 64 esempi di apprendimento e 96 nuovi esempi di verifica
per seed, il nucleo restituisce il candidato corretto in $571/576$ casi.
Il Transformer online con replay ottiene $570/576$, mentre il controllo
episodico con distanza fissa ottiene $574/576$. Il test esatto appaiato per seed
ha $p=1.00$: non c'è una dimostrazione di vantaggio. La perdita logaritmica
favorisce il Transformer, $0.1681$ contro $0.2037$.

Il nucleo accetta 533 risposte, tutte corrette nell'esperimento, e richiede
conferma negli altri 43 casi. La generalizzazione sulle classi precedenti
non peggiora nei sei flussi; non segue una garanzia fuori da tali flussi.
I tre canali insieme forniscono circa il 99\% di accuratezza, i singoli
canali circa il 25\%, due canali circa il 49\%. Il risultato riguarda colori,
toni e indizi testuali artificiali, non intelligenza generale.

La curvatura e le equazioni native del vecchio motore MGD non determinano
queste predizioni: il componente aggiunto è un classificatore episodico con
metrica appresa. Attribuire il risultato a una superiorità della teoria MGD
sarebbe un errore di inferenza. Tempi CPU non equivalgono a energia; non sono
stati misurati joule sul telefono.

\begin{thebibliography}{9}
\bibitem{nca} Goldberger et al., \emph{Neighbourhood Components Analysis},
NeurIPS 2004. \url{https://www.cs.utoronto.ca/~hinton/absps/nca.pdf}
\bibitem{cls} \emph{Organizing memories for generalization in complementary
learning systems}, Nature Neuroscience, 2023.
\url{https://www.nature.com/articles/s41593-023-01382-9}
\bibitem{der} Buzzega et al., \emph{Dark Experience for General Continual Learning},
2020. \url{https://arxiv.org/abs/2004.07211}
\bibitem{mamba} Lahoti et al., \emph{Mamba-3}, 2026.
\url{https://arxiv.org/abs/2603.15569}
\bibitem{bind} Girdhar et al., \emph{ImageBind}, 2023.
\url{https://arxiv.org/abs/2305.05665}
\bibitem{nl} Behrouz et al., \emph{Nested Learning}, 2025.
\url{https://abehrouz.github.io/files/NL.pdf}
\end{thebibliography}
\end{document}
